\documentclass[10pt,a4paper]{amsart}
\usepackage[margin=19mm]{geometry}
\usepackage{amsmath,amssymb,mathtools}
\usepackage[scr=rsfs,cal=euler]{mathalfa}
\usepackage{xfrac}
\usepackage[unicode,hidelinks,bookmarks=false]{hyperref}
\newcommand{\ie}{i.e.\ }
\newcommand{\cf}{cf.\ }
\newcommand{\resp}{respectively\ }
\newcommand{\Subsection}{\subsection}
\providecommand{\nobreakdash}{\nobreak\hbox{-}}
\setlength{\emergencystretch}{2em}
\multlinegap0pt
\usepackage{dsfont}
\usepackage{amssymb}
\usepackage{csquotes}
\usepackage{mathtools}
\usepackage{braket}
\usepackage[shortlabels]{enumitem}
\usepackage{xfrac}
\usepackage{tikz}
\newtheorem{lemma}{Lemma}
\newcommand{\R}{\mathbb{R}}
\title{Continuity equation, Sobolev-to-Lipschitz and Ricci curvature bound for the metric measure space of probability measures}
\author{Alessandro Pinzi}
\date{Source: arXiv:2511.21204v2 (CC BY 4.0)}
\begin{document}
\maketitle

\noindent\emph{Source and licence.} Continuity equation, Sobolev-to-Lipschitz and Ricci curvature bound for the metric measure space of probability measures. Version arXiv:2511.21204v2 (CC BY 4.0), distributed by arXiv under the Creative Commons Attribution 4.0 International licence, http://creativecommons.org/licenses/by/4.0/, as recorded in the arXiv licence metadata for this exact version and shown on the abstract page https://arxiv.org/abs/2511.21204v2. Full licence notice: \texttt{licenses/arxiv-2511.21204-CC-BY-4.0.txt}.\par\smallskip

\noindent\textit{Editorial scope.} Counted unit: the article's lemma appendix (source0.tex lines 2800--2803). The article's complete original proof of it is delivered with it (source0.tex lines 2805--2897, one \texttt{proof} environment of 93 lines). No mathematics is written, repaired or re-derived: the statement and the proof are the article's own lines, in the article's own notation, and no retained line is substituted or re-flowed.  No further source-local context is needed: every cross-reference of every retained line resolves inside the two delivered spans.  Nothing was omitted from any retained span, so the package carries no omission record.  No retained line contains a citation, so the wrapper carries no bibliography and no bibliography entry.  No retained line contains a figure environment, an image-inclusion command or drawing code, so no image is delivered and the package declares no asset dependency.  Editorial formatting only, in addition: an AMS \texttt{amsart} preamble that restates the article's own declarations and macros used by the retained lines, namely the environments \texttt{lemma} on separate counters, together with the standard packages those lines need.  The retained environments print in this document as Lemma 1 in order of appearance, whereas the article prints the same items with the numbering of its own full document.  The complete native source of the article as distributed in the arXiv e-print for this exact version is bundled unchanged as \texttt{sources/189836/source0.tex}.
\begin{lemma}\label{lemma: capacity appendix}
    Assume that $\nu \ll \mathcal{L}^d$. Then for all $p\geq 1$
    \[\operatorname{cap}_{p,\nu}(\Delta) = \widetilde{\operatorname{cap}}_{p,\nu}(\Delta).\]
\end{lemma}

\begin{proof}
Clearly, $\operatorname{cap}_{p,\nu}(\Delta) \geq \widetilde{\operatorname{cap}}_{p,\nu}(\Delta)$, so we focus on the other inequality.
    Let us call $\mathcal{C}(h) := \int |h^p| + |\nabla h|^p d\nu\otimes \nu$. We are going to prove that for any $h\in C_b^1(\R^{2d})$ satisfying $h\geq 1$ on $\Delta$ and for any $\delta>0$, there exists a function $g\in C_b^1(\R^{2d})$ such that $g\leq 1$ everywhere, $g=1$ on $\Delta$ and 
    \begin{equation}\label{eq: proof of cap lemma}
       \mathcal{C}(g) \leq \mathcal{C}(h) + \delta.
    \end{equation}

    \textbf{Step 1}: define $\tilde{h}(x,y) := h(x,y)\wedge 1$. Notice that $\tilde{h}$ is a Lipschitz function, and satisfies $\tilde{h}\leq 1$, $\tilde{h} = 1$ on $\Delta$ and $\mathcal{C}(\tilde{h})\leq \mathcal{C}(h)$, thanks to the locality of the gradient for Lipschitz functions, i.e. 
    \begin{equation}\label{eq: gradient of truncated h}
    \nabla \tilde{h}(x,y) = \nabla h(x,y) \mathds{1}_{\{h\leq 1\}}(x,y) \quad \mathcal{L}^{2d}\text{-a.e.},
    \end{equation}
    and the fact that $\nu \ll \mathcal{L}^d$.

    \textbf{Step 2}: we claim that there exists a Lipschitz function $\overline{h}$ such that $\overline{h}\leq 1$, $\mathcal{C}(\overline{h})\leq \mathcal{C}(\tilde{h}) + \frac{\delta}{2}$ and $\overline{h} = 1$ on $\Delta_{\varepsilon}:= \Delta+B(0,\varepsilon)$, for some $\varepsilon>0$, i.e. there exists a `strip' of length $\varepsilon$ around the diagonal $\Delta$ where the function $\overline{h}$ is constantly $1$.
    \\
    To this aim, define $\pi_\Delta$ and $\pi_{\Delta^\perp}$, respectively, the orthogonal projection over $\Delta$ and $\Delta^\perp$, i.e. 
    \begin{align*}
        \pi_\Delta(x,y) = \left(\frac{x+y}{2}, \frac{x+y}{2}\right), \quad \pi_{\Delta^\perp}(x,y) = \left( \frac{x-y}{2}, - \frac{x-y}{2} \right), \\
        (x,y) = \pi_\Delta(x,y) + \pi_{\Delta^\perp}(x,y), \quad \|(x,y)\|^2 = \|\pi_\Delta(x,y)\|^2 + \|\pi_{\Delta^\perp}(x,y)\|^2.
    \end{align*}
    For any $\varepsilon>0$, define $p_\varepsilon:\R^{2d}\to \R^{2d}$ as follows 
    \begin{equation}
        p_\varepsilon(x,y) := \begin{cases}
            \pi_{\Delta}(x,y) \quad & \text{ if } (x,y) \in \Delta_\varepsilon
            \\
            (x,y) - \varepsilon \frac{\pi_{\Delta^\perp}(x,y)}{\|\pi_{\Delta^\perp}(x,y)\|} & \text{ otherwise}.
        \end{cases}
    \end{equation}
    In the next step, we will prove that $p_\varepsilon$ is $1$-Lipschitz for any $\varepsilon>0$. For now, assume that is true and we continue defining $h_\varepsilon(x,y) := \tilde{h}(p_\varepsilon(x,y))$. 
    %Thanks to \eqref{eq: gradient of truncated h} and the $1$-Lipschitzianity of $p_\varepsilon$, we have that 
    %\[|\nabla \overline{h}(x,y)| \leq |\nabla \tilde{h}(p_{\varepsilon}(x,y))|  = |\nabla h(p_{\varepsilon}(x,y))|\mathds{1}_{\{h\circ p_\varepsilon \leq 1\}} \quad  \mathcal{L}^{2d}-a.e.,\]
    %and in particular the right hand side is upper-semicontinuous. 
    We conclude the proof of this step if 
    \[\limsup_{\varepsilon\searrow 0} \mathcal{C}(h_\varepsilon)\leq \mathcal{C}(\tilde{h}).\]
    To this aim, notice that $h_\varepsilon \to \tilde{h}$ pointwise as $\varepsilon \searrow 0$, so that 
    \[\int |h_\varepsilon|^p d\nu\otimes \nu \to \int|\tilde{h}|^p d\nu\otimes \nu\]
    by dominated convergence theorem. Regarding the gradients, notice that $|\nabla \tilde{h}(x,y)|=  |\nabla h(x,y)| \mathds{1}_{\{ h\leq 1\} }(x,y)$ almost everywhere, and the right hand side is an upper semicontinuous function. Then, thanks to the straightforward fact that $(p_\varepsilon)_\#(\nu\otimes \nu) \to \nu\otimes \nu$ in the narrow convergence and 
    \[|\nabla h_\varepsilon(x,y)| \leq |\nabla \tilde{h}(p_\varepsilon(x,y))| = |\nabla h(p_\varepsilon(x,y))| \mathds{1}_{\{h\leq 1\}}(p_\varepsilon(x,y)) \quad \mathcal{L}^{2d}\text{-a.e.},\] 
    we have that 
    \begin{align*}
        \limsup_{\varepsilon\searrow 0} &  \int |\nabla h_\varepsilon|^p(x,y) d\nu\otimes \nu(x,y) \leq \limsup_{\varepsilon\searrow 0} \int |\nabla h(x,y)|^p\mathds{1}_{ \{h\leq 1\} } (x,y ) d\big((p_\varepsilon)_\#(\nu \otimes \nu)\big)(x,y) 
        \\
        \leq & 
        \int |\nabla h(x,y)|^p\mathds{1}_{\{h\leq 1\}}(x,y) d\nu\otimes \nu(x,y) = \int |\nabla \tilde{h}(x,y)|^p d\nu\otimes \nu(x,y).
    \end{align*}

    \textbf{Step 3}: proof that $p_\varepsilon$ is $1$-Lipschitz for all $\varepsilon>0$. We need to verify that $\|p_\varepsilon(x,y) - p_\varepsilon(z,w)\| \leq \|(x,y) - (z,w)\|$ for all $(x,y),(z,w)\in \R^d\times \R^d$, where $\|\cdot\|$ is the Euclidean norm of $\R^{2d}$. There are three cases:
    \begin{itemize}
        \item $(x,y),(z,w) \in \Delta_\varepsilon$, then it is straightforward because any orthogonal projection is $1$-Lipschitz;
        \item $(x,y)\in \Delta_\varepsilon$ and $(z,w)\notin \Delta_\varepsilon$, then notice that $\|\pi_{\Delta^\perp}(x,y)\|<\varepsilon$ and $\|\pi_{\Delta^\perp}(z,w)\|\geq \varepsilon$. Then
        \begin{align*}
            \|p_\varepsilon(x,y) - p_\varepsilon(z,w)\|^2 
             = &
            \| \pi_\Delta(x,y) - \pi_\Delta(z,w) -\left(1- \frac{\varepsilon}{\|\pi_{\Delta^\perp}(z,w)\|} \right)\pi_{\Delta^\perp}(z,w) \|^2
            \\
             = &
            \|\pi_\Delta(x,y) - \pi_\Delta(z,w)\|^2 + \left(\|\pi_{\Delta^\perp}(z,w)\| - \varepsilon\right)^2
            \\
            (\star) \leq &
            \|\pi_\Delta(x,y) - \pi_\Delta(z,w)\|^2 + \|\pi_{\Delta^\perp}(x,y) - \pi_{\Delta^\perp}(z,w)\|^2 
            \\
            = & \|(x,y) - (z,w)\|^2,
        \end{align*}
        where in $(\star)$ we used that $\|\pi_{\Delta^\perp}(z,w)\| - \varepsilon\leq \|\pi_{\Delta^\perp}(z,w)\| - \|\pi_{\Delta^\perp}(x,y)\|\leq \|\pi_{\Delta^\perp}(x,y) - \pi_{\Delta^\perp}(z,w)\|$;
        \item $(x,y),(z,w) \notin \Delta_\varepsilon$. Then 
        \begin{equation}\label{eq: outside strip}
        \begin{aligned}
            \|p_\varepsilon(x,y) & - p_\varepsilon(z,w)\|^2 
            = 
            \|\pi_{\Delta}(x,y) - \pi_{\Delta}(z,w)\|^2
            \\
            & + \left\|\left(1 - \frac{\varepsilon}{\|\pi_{\Delta^\perp}(x,y) \|}\right)\pi_{\Delta^\perp}(x,y) - \left(1 - \frac{\varepsilon}{\|\pi_{\Delta^\perp}(z,w) \|}\right)\pi_{\Delta^\perp}(z,w)\right\|^2.
        \end{aligned}
        \end{equation}
        For simplicity, call $v_1 =\pi_{\Delta^\perp}(x,y)$ and $v_2 =\pi_{\Delta^\perp}(z,w)$, together with $\alpha_i := 1- \varepsilon/\|v_i\|$, for $i=1,2$. Then
        \begin{align*}
            \|\alpha_1 v_1 - \alpha_2 v_2\| = & 
            \|(\alpha_1-\alpha_2)v_1 + \alpha_2 (v_1 - v_2)\|
            \leq 
            |\alpha_1 - \alpha_2|\|v_1\| + \alpha_2 \|v_1 - v_2\|
        \end{align*}
        and 
        \begin{align*}
            |\alpha_1 - \alpha_2|\|v_1\| = \varepsilon\left|\frac{1}{\|v_2\|} - \frac{1}{\|v_1\|}\right|\|v_1\|\leq\varepsilon \frac{\|v_1-v_2\|}{\|v_2\|},
        \end{align*}
        so that 
        \[\|\alpha_1 v_1 - \alpha_2 v_2\| \leq \frac{\varepsilon}{\|v_2\|}\|v_1 - v_2\| + \|v_1-v_2\| - \frac{\varepsilon}{\|v_2\|}\|v_1-v_2\| = \|v_1-v_2\|.\]
        Substituting in \eqref{eq: outside strip} we conclude that $p_\varepsilon$ is $1$-Lipschitz.
    \end{itemize}

    \textbf{Step 4}: convolution and conclusion. Consider $\overline{h}$ as in step 2. Then, for any $0 <\alpha < \varepsilon$, define the functions $g_\alpha:= \overline{h} * \rho_\alpha$, where $\rho_\alpha(x,y) := \frac{1}{\alpha^{2d}}\rho\left(\frac{(x,y)}{\alpha}\right)$ and $\rho$ is a mollifier with support in $B(0,1)$. Notice that each $g_\alpha\in C_b^1(\R^{2d})$ satisfies $h\leq 1$ and $h = 1$ on $\Delta$. We conclude observing that, up to selecting a suitable subsequence $\alpha_n \to 0$, $|g_\alpha|^p + |\nabla g_\alpha|^p \to |\overline{h}|^p + |\nabla \overline{h}|^p$ pointwise almost everywhere, and then by dominated convergence theorem we have that 
    \[\lim_{\alpha \to 0}\mathcal{C}(g_\alpha) = \mathcal{C}(\overline{h}). \qedhere \]
\end{proof}

\end{document}
